% !PRINTED 205
\weberpage{205}

%<<CONT
\[
\frac{1-x^2}{1-\varkappa^2 x^2}, \qquad \frac{1-\varkappa^2 x^2}{\varkappa^2(1-x^2)}, \qquad \frac{1}{\varkappa^2 x^2}
\]

die Wurzeln von $B = 0$, $C = 0$, $D = 0$. Die Bedeutung dieser
letzteren Wurzel ist aber % [sic] The source prints the singular although three roots precede.

\[
\operatorname{sn}\left(\frac{(4\mu+1)K + 4\mu' iK'}{n}\right)^2,
\]
\[
\operatorname{sn}\left(\frac{(4\mu+1)K + (4\mu'+1)iK'}{n}\right)^2,
\]
\[
\operatorname{sn}\left(\frac{4\mu K + (4\mu'+1)iK'}{n}\right)^2,
\]

und hiernach sind durch Auflösung der Gleichung $A = 0$ alle
Größen von der Form

\[
\operatorname{sn}\left(\frac{\mu K + \mu' iK'}{n}\right)^2,
\]

worin $\mu$, $\mu'$ beliebige ganze Zahlen sind, \gesperrt{rational bestimmt}.

\begin{center}
\S\ 62. \quad \gesperrt{Die Abelschen Relationen.}
\end{center}

Für eine nähere Untersuchung der algebraischen Natur der
Periodenteilungsgleichung ist ein System von Relationen zwischen
ihren Wurzeln von großer Wichtigkeit, zu dessen Ableitung wir
jetzt übergehen \textsuperscript{1}).

Wir betrachten die Summe:

\weq{1}{\sum^{\nu} e^{\frac{8\nu\nu'\pi i}{n}} \frac{\vartheta_{11}\left(u + \frac{2\nu}{n}\right)}{\vartheta_{g_1,g_2}\left(u + \frac{2\nu}{n}\right)},}

genommen nach $\nu$ über ein volles Restsystem für den (ungeraden)
Modul $n$. $\nu'$ ist eine beliebige ganze Zahl und $g_1$, $g_2$ eine der
drei geraden Charakteristiken $(0,0)$, $(0,1)$, $(1,0)$. Diese Summe
ist unabhängig von dem besonderen Restsystem, das man für $\nu$
genommen hat.

\begingroup
\footnotesize
\textsuperscript{1}) Diese Relationen rühren von \gesperrt{Abel} her (Oeuvres complètes, Bd. I,
S. 523; Bd. II, S. 251 der neuen Ausgabe). Der oben gegebene Beweis
dieser Relationen schließt sich an \gesperrt{Hermite} an (Crelles Journal, Bd. 32,
S. 283). Zu erwähnen ist noch \gesperrt{Sylow}, Christiania Videnskabsselskabs
Forhandlinger 1864 und 1871. \gesperrt{Kronecker}, Berichte der Berliner Akademie,
19. Juli 1875. \gesperrt{Engel}, Berichte der Sächsischen Gesellschaft der Wissenschaften, 31. Juli 1884.
\endgroup
